\section{Discussion}
\label{sec:discussion}

\subsection{Claims}
\label{sec:claims}

\begin{description}[leftmargin=2.2em,style=nextline]
  \item[\textbf{C1 (H2, headline).}]
    Because the model \emph{is} the prior, model error enters the analysis
    unbuffered, and its cost is not a constant offset: for strong-constraint
    4D-Var it \emph{collapses the marginal value of observations} ($6.4\times$;
    $7.0\times$ at the benchmark inflation). The ensemble filters keep most of it
    ($1.1\times$; $1.6$--$1.7\times$ at the benchmark inflation).
    \caveat{Established for strong-constraint 4D-Var; the filters' collapse
    reported in earlier drafts ($1.9\times$) came from a mis-specified \Sz{} and
    does not survive the fix (Section~\ref{sec:marginal}). Lorenz-63
    shows weak-constraint 4D-Var recovers much of the loss
    (Section~\ref{sec:weak4dvar}); the corresponding Lorenz-96 measurement on the
    density grid does not yet exist, so C1 is stated for strong-constraint
    4D-Var and not for the classical framework as a whole.}
  \item[\textbf{C2 (H3a).}]
    The state representation fixes the conditioning of the inverse problem.
    Identical physics in a different state variable produces large, predictable
    skill differences via $q\approx\nabla^2\psi$ and its $K^2$ amplification.
  \item[\textbf{C3 (H1).}]
    Markovian factorisation forecloses \emph{learning} the observation
    configuration, and buys robustness to it for free. DA is the scheme least
    affected by a change of observing system ($\times 1.06$--$1.20$); an amortised
    estimator trained on one configuration is not ($\times 2.7$--$3.3$). Trained on
    the \emph{distribution} of configurations, it loses little
    ($\times 1.24$--$1.32$), matches the fixed-grid model on that grid, and beats DA
    from about 20 observations per window at \Sz{} and at every density at \So{}
    (Section~\ref{sec:ordering}).
  \item[\textbf{C4 (H1 $+$ H3a $\Rightarrow$ S-b) --- \emph{not supported as
    originally stated}.}]
    The intended claim was that the ensemble-Kalman class, carrying
    $\PsiNG \equiv 0$ by construction, has a structurally impoverished posterior.
    \textbf{Neither dispersion nor shape measurements support it}
    (Section~\ref{sec:nongaussian}). ETKF and EnKF are the best-calibrated
    schemes measured (spread/RMSE $0.98$, $0.97$ against a target of $0.967$),
    while every scheme with a full non-Gaussian branch is under-dispersed by
    $1.4\times$ to $2.3\times$. Rank histograms --- which could see a shape
    deficit that spread/RMSE cannot --- show the filters' residual shape error,
    once spread and mean bias are corrected, in the same range as
    PredictStateCFM's and below SDA's. \textbf{C4 is therefore not supported as a
    discriminating claim}, and S-b rests on the accuracy/dispersion trade being
    unresolved by \emph{every} class rather than on a deficit specific to one.
  \item[\textbf{C5 (attribution).}]
    Each hypothesis can be relaxed separately and the resulting change
    attributed to a named term of~\eqref{eq:partition}. This is what makes the
    paper a diagnosis rather than a benchmark.
  \item[\textbf{C6 (H2 $\times$ H3b).}]
    The representation's sensitivity to $(\thv,\zv)$ is \emph{not a proxy} for
    the sensitivity of $p(\xv_1\mid\Cv)$ to $(\thv,\zv)$. DA conflates them
    because its estimator is the model; a conditioned amortised estimator
    separates them; and the gap is large --- $2.0$--$2.2\times$ against $0$--$2\%$
    (Section~\ref{sec:sensitivity}).
\end{description}

\subsection{Limitations and threats to validity}
\label{sec:limitations}

\paragraph{The \Sz{}/\So{} axis means different things per class.}
\So{} is a \emph{no-op} for model-free learned schemes; their $\So/\Sz\approx 1.00$
carries no information and is never tabled as robustness
(Table~\ref{tab:classes}). Even for conditioned schemes, being handed biased
parameters is not the same exposure as \emph{integrating} a biased model for 500
steps: one is an input perturbation, the other compounds. Section~\ref{sec:marginal}'s
intra-classical result needs neither caveat.

\paragraph{``Your baselines were under-tuned.''}
The standard reply to any DA-versus-learned comparison. Our mitigation is
evidential: localisation and inflation sensitivity studies found the QG
localisation radius untuned \emph{even at the reference density}, with an optimum
that moves with observation density --- itself an instance of what H1 forecloses
--- and a 4D-Var variance-scaling sweep returned a \emph{negative} result, the
existing defaults already being best. On Lorenz-96 the objection was partly
right, and we disclose it: the DA forward model had not received the
fast-variable coupling weights and \Sz{} ran at the \So{}-tuned inflation; fixing
both lowered ETKF's \Sz{} RMSE by 21\%. Every number here is the fixed one.
Section~\ref{sec:weak4dvar} is the stronger answer: the classical family is
represented by its best member, not its most convenient one.

\paragraph{Uneven scheme coverage.}
Only Lorenz-96 carries all five classes (Table~\ref{tab:coverage}). Cross-class
claims are made there; the other two systems carry specific arguments. This is a
design choice, but it does mean the generality of the cross-class ordering rests
on one system.

\paragraph{Different configurations must not be cross-quoted.}
Tables~\ref{tab:price} and~\ref{tab:marginal} come from different configurations
--- different observation counts, different randomisation breadth. Ratios from
one are not comparable to levels in the other. The same holds across test sets
(regular grid versus random observing system) and across training protocols
(P1 versus benchmark default, 400 versus 1200 epochs): the learned ranking itself
changes between them.

\paragraph{The circularity objection.}
An amortised estimator needs a simulator to generate its training data. If that
simulator is the truth model, the comparison has handed it the truth --- and in
operations there is no truth model. This is the deepest objection to this line of
work and we do not claim to answer it fully. The framing is therefore about
\emph{where prior information enters and what it costs}, not ``learned beats
physical''.
\todo{A self-supervised objective --- a weak-constraint variational cost
evaluated on the estimator's own output, with no ground-truth supervision --- is
implemented and is a direct, if partial, answer. Not yet run; decide whether it
belongs here or in the companion paper.}

\paragraph{Idealised testbeds.}
All three systems are idealised. That is what makes ``matched everything''
definable precisely enough for a marginal-value measurement to mean anything,
but transfer to operational systems is an open question we do not test.

\subsection{Positioning}
\label{sec:positioning}

\todo{Every citation is from memory and must be verified before a bibliography
exists. Nothing in \texttt{refs.bib} has been checked, and no \texttt{$\backslash$cite}
commands are used yet.}

The DA review literature frames the field's standard commitments; the
assimilation-in-the-unstable-subspace line gives the closest existing account of
\emph{why} DA's effective degrees of freedom are limited; weak-constraint 4D-Var
is the canonical relaxation of H2; innovation-based diagnostics supply the
standard consistency checks, whose assumptions are violated when the estimator is
trained to fit the observations. On the learned side, the
DA-with-learned-model-error-correction line and the score-based and flow-based
assimilation line are the nearest neighbours.

\paragraph{The gap.}
The model-error literature asks how to \emph{correct} the model; the
unstable-subspace literature explains the dimensionality of the analysis; the
learned-DA literature proposes hybrids. As far as our survey goes, none of them
\emph{measures the marginal value of observations as a function of model error},
or attributes the deficit to named components of the inference operator.
Section~\ref{sec:marginal} is, on that reading, an unreported effect.

\subsection{Open experiments}
\label{sec:open}

\todo{This list is provisional. It was drawn up against the previous,
hypothesis-ordered outline; it should be re-derived from the structure above now
that the results lead with the cross-system overview, and pruned to what the
paper actually needs rather than everything that could be measured.}

\begin{description}[leftmargin=2.2em,style=nextline]
  \item[\textbf{Weak-constraint 4D-Var on Lorenz-96.}] Bounds C1. Implemented;
    the benchmark run does not exist. Highest value, lowest cost.
  \item[\textbf{H1 isolated.}] Filter vs.\ window-smoother vs.\ window-amortised
    at matched model, observations and compute --- what turns
    Section~\ref{sec:ordering}'s ordering into evidence.
  \item[\textbf{Model-error dose--response.}] Sweep model error continuously
    rather than the binary \Sz{}/\So{}, giving the \emph{same} biased parameters
    to the classical forward model and to the learned conditioning inputs.
  \item[\textbf{The marginal-value surface.}] Observation density $\times$ model
    error for every class on matched axes. A factorial grid of random observing
    systems (6 observation counts $\times$ 4 fast-channel counts, the same layouts
    for every class) now covers the density axes at \Sz{} and \So{}; the missing
    part is a continuous model-error axis.
  \item[\textbf{The two sensitivities on one grid.}] Model sensitivity and
    posterior sensitivity over observation density $\times$ perturbation
    amplitude; the crossover density is the paper's second deliverable number
    after Section~\ref{sec:marginal}'s collapse factor.
\end{description}
